For a consequential transition
\[
x_t \xrightarrow{a_t} x_{t+1},
\]
admission requires explicit obligations
\[
O = \{\text{identity}, \text{authority}, \text{evidence},
\text{ACP binding}, \text{XAL verification}, \text{continuity},
\text{effect scope}\}.
\]
If any mandatory obligation is missing,
\[
\neg\bigwedge_{o \in O} o \implies
\operatorname{Disposition} \in \{\text{HOLD}, \text{REJECT}\}.
\]

\begin{definition}[Consequential transition]
A transition $x_t \xrightarrow{a_t} x_{t+1}$ is \emph{consequential} if it
changes persistent state, invokes a tool with external effect, modifies
code or configuration, or affects another principal.
\end{definition}

\begin{definition}[Obligation discharge]
An obligation $o \in O$ is \emph{discharged} when a content-bound reference
exists that an independent verifier can recompute: a digest, a lease, a
signed program, or a receipt binding the obligation to the specific
transition.
\end{definition}

The estate reference implementation encodes these obligations as a typed
set. A transition request carries optional fields for each obligation;
the gate derives the missing set and returns REJECT when any mandatory
obligation is absent, HOLD when evidence is present but not content-bound,
and ADMIT only when every obligation is discharged with content-addressed
evidence.

\begin{theorem}[Reject absorption]\label{thm:reject-absorbing}
Once a hard fault exists, downstream components cannot convert the
transition into ADMIT, PASS, or PROMOTE.
\end{theorem}

\begin{proof}
By the definition of the disposition function: a hard fault forces the
disposition to REJECT at the transition gate, and every downstream stage
(compilation, execution, replay, promotion) consumes the disposition
without the authority to override it. Formally,
$\texttt{forcesReject}(d) \iff d = \texttt{REJECT}$, and
$\texttt{hard\_fault\_not\_admit}$ establishes $d \neq \texttt{ADMIT}$
whenever the fault exists.
\end{proof}

The Lean formalization \texttt{RejectAbsorption.lean} proves both
\texttt{reject\_is\_absorbing} and \texttt{hard\_fault\_not\_admit}
structurally.
